\documentclass[11pt]{article}

\usepackage[margin=1in]{geometry}
\usepackage{amsmath, amssymb, booktabs, longtable}
\usepackage{natbib}
\usepackage{hyperref}

\title{Methodological Notes for Event Tests on Daily CARDI}
\author{CARDI Project}
\date{May 19, 2026}

\begin{document}
\maketitle

\section{Purpose}

This note documents the econometric methods implemented in
\texttt{Code/Plot\_Code/Event\_test.R}. The goal is to test whether
carbon-related regulatory events are associated with structural changes in the
daily CARDI time series. The script uses three daily CARDI series:
\[
\text{CARDI}_{1P,t}, \qquad \text{CARDI}_{5P,t}, \qquad \text{CARDI}_{10P,t}.
\]

The script implements four empirical components:
\begin{enumerate}
  \item event clustering and overlap diagnostics;
  \item known-breakpoint Chow tests;
  \item interrupted time-series regressions with Newey--West standard errors;
  \item two-regime Markov-switching regime analysis.
\end{enumerate}

It also preserves a reference-style event-dummy regression based on smoothed
CARDI log changes, with the important modification that event-dummy coefficients
are reported using Newey--West standard errors.

\section{Data}

Let \( y_t^{(q)} \) denote the CARDI series for quantile level
\(q \in \{1P,5P,10P\}\). The daily CARDI data are loaded from:
\[
\texttt{Data/Processed/FRM\_Carbon\_risk.csv}.
\]

The event data are loaded from:
\[
\texttt{Data/raw/Important\_Carbon\_Events.xlsx}.
\]

Let \( \mathcal{E} = \{e_1,\ldots,e_J\} \) denote the set of event dates.
If an event occurs on a non-trading day, the script maps the event to the
nearest observed CARDI trading date.

\section{Event Clustering}

Several policy events occur close to one another. Treating nearby events as
independent shocks can make local event-window tests difficult to interpret.
The script therefore constructs event clusters before running the main tests.

Let the sorted event dates be:
\[
e_1 \leq e_2 \leq \cdots \leq e_J.
\]
For a configurable clustering threshold \(G\), currently set to 30 calendar
days, event \(e_j\) belongs to the same cluster as \(e_{j-1}\) if
\[
e_j - e_{j-1} \leq G.
\]
Otherwise, \(e_j\) starts a new cluster.

For each cluster \(c\), the script records:
\[
\left(
\text{cluster\_id}_c,\;
\text{start}_c,\;
\text{end}_c,\;
\text{representative\_date}_c,\;
N_c
\right),
\]
where \(N_c\) is the number of events in the cluster. The representative event
date is the first event date in the cluster.

\section{Known-Breakpoint Chow Test}

\subsection{Motivation}

The Chow test evaluates whether regression coefficients are stable across a
known breakpoint. The test was introduced by \citet{chow1960}. In this
application, each event date is treated as a known candidate breakpoint.

\subsection{Local Event Window}

For an event date \(e\) and a symmetric trading-day window \(h\), the local
sample is:
\[
\mathcal{T}(e,h)
= \{t: -h \leq \tau_t(e) \leq h\},
\]
where \(\tau_t(e)\) is the trading-day distance between date \(t\) and the
nearest CARDI trading date for event \(e\).

\subsection{Restricted and Unrestricted Models}

For each CARDI variable \(y_t^{(q)}\), the restricted model is:
\begin{equation}
y_t^{(q)}
= \alpha + \beta t + u_t.
\end{equation}

The unrestricted model allows both the intercept and slope to change after the
event:
\begin{equation}
y_t^{(q)}
= \alpha + \beta t
+ \gamma D_t(e)
+ \delta \left[t \times D_t(e)\right]
+ u_t,
\end{equation}
where:
\[
D_t(e) =
\begin{cases}
1, & t \geq e,\\
0, & t < e.
\end{cases}
\]

\subsection{Test Statistic}

Let \(RSS_R\) be the residual sum of squares from the restricted model and
\(RSS_U\) be the residual sum of squares from the unrestricted model. The Chow
statistic is:
\begin{equation}
F
=
\frac{(RSS_R - RSS_U)/q}{RSS_U/(n-k_U)},
\end{equation}
where \(q\) is the number of restrictions, \(n\) is the local sample size, and
\(k_U\) is the number of parameters in the unrestricted model.

The null hypothesis is:
\[
H_0: \gamma = 0 \quad \text{and} \quad \delta = 0.
\]

Rejecting \(H_0\) indicates evidence of a known structural break around the
event date.

\section{Interrupted Time-Series Regressions}

\subsection{Motivation}

Interrupted time-series models are commonly used to estimate whether an
intervention is associated with an immediate change, a persistent level shift,
or a trend change. See \citet{bernal2017} for a tutorial treatment of segmented
regression designs for intervention analysis.

\subsection{Regression Specification}

For each CARDI variable \(y_t^{(q)}\), event \(e\), and local window
\(\mathcal{T}(e,h)\), the script estimates:
\begin{equation}
y_t^{(q)}
= \alpha
+ \beta_1 \text{Pulse}_{t,e}
+ \beta_2 \text{Post}_{t,e}
+ \beta_3 \text{Trend}_{t,e}
+ \varepsilon_t.
\end{equation}

The event variables are defined as:
\begin{align}
\text{Pulse}_{t,e}
&=
\mathbf{1}\{t=e\},\\
\text{Post}_{t,e}
&=
\mathbf{1}\{t>e\},\\
\text{Trend}_{t,e}
&=
\max\{\tau_t(e),0\}.
\end{align}

The coefficients have the following interpretation:
\begin{itemize}
  \item \(\beta_1\): same-day event shock;
  \item \(\beta_2\): persistent post-event level shift;
  \item \(\beta_3\): post-event trend change.
\end{itemize}

\subsection{Newey--West Standard Errors}

Because daily CARDI is a time series, residuals may be serially correlated and
heteroskedastic. The script therefore reports Newey--West heteroskedasticity
and autocorrelation consistent standard errors \citep{neweywest1987}.

Let \(X\) be the matrix of regressors and \(\hat{\varepsilon}_t\) be the OLS
residuals. The Newey--West covariance estimator can be written as:
\begin{equation}
\widehat{\mathrm{Var}}_{\mathrm{NW}}(\hat{\beta})
=
(X'X)^{-1}
\left(
\sum_{\ell=-L}^{L} w_\ell
\sum_{t=|\ell|+1}^{T}
x_t x_{t-|\ell|}' \hat{\varepsilon}_t \hat{\varepsilon}_{t-|\ell|}
\right)
(X'X)^{-1},
\end{equation}
where \(L\) is the truncation lag and the Bartlett weights are:
\begin{equation}
w_\ell = 1 - \frac{|\ell|}{L+1}.
\end{equation}

The script uses:
\[
L=5.
\]

\section{Reference-Style Event-Dummy Regression}

\subsection{Smoothed CARDI Series}

The script preserves the legacy plotting and regression logic from the earlier
\texttt{Event\_test.R}. First, each CARDI series is smoothed using LOESS over
the daily time index:
\begin{equation}
\widetilde{y}_t^{(q)}
=
\mathrm{LOESS}\left(y_t^{(q)}; \mathrm{span}=0.1\right).
\end{equation}

The dependent variable is the log change in smoothed CARDI:
\begin{equation}
\Delta \widetilde{y}_t^{(q)}
=
100 \times
\left[
\log\left(\widetilde{y}_t^{(q)}\right)
-
\log\left(\widetilde{y}_{t-1}^{(q)}\right)
\right].
\end{equation}

\subsection{Trailing Event Dummy}

For each lookback window \(m \in \{7,30,60\}\), define:
\begin{equation}
\text{DummyEvent}_{t,m}
=
\mathbf{1}
\left\{
\exists e_j \in \mathcal{E}
\text{ such that }
t-(m-1) \leq e_j \leq t
\right\}.
\end{equation}

The regression is:
\begin{equation}
\Delta \widetilde{y}_t^{(q)}
=
\alpha
+ \theta_m \text{DummyEvent}_{t,m}
+ \eta_t.
\end{equation}

The coefficient \(\theta_m\) measures whether smoothed CARDI log changes are
different on days following recent carbon-related events. The script reports
Newey--West standard errors for \(\theta_m\).

\section{Markov-Switching Regime Analysis}

\subsection{Motivation}

The Markov-switching model tests whether CARDI alternates between low-risk and
high-risk latent regimes. The approach follows the regime-switching framework
introduced by \citet{hamilton1989}.

\subsection{Two-Regime Model}

For each CARDI series, the script estimates a two-regime model:
\begin{equation}
y_t^{(q)}
=
\mu_{s_t}^{(q)} + \varepsilon_t,
\qquad
\varepsilon_t \sim N\left(0,\sigma_{s_t}^2\right),
\end{equation}
where \(s_t \in \{1,2\}\) is an unobserved Markov state.

The state evolves according to:
\begin{equation}
\Pr(s_t=j \mid s_{t-1}=i)
=
p_{ij},
\qquad i,j \in \{1,2\}.
\end{equation}

The transition matrix is:
\begin{equation}
P =
\begin{pmatrix}
p_{11} & p_{12}\\
p_{21} & p_{22}
\end{pmatrix}.
\end{equation}

\subsection{High-Risk Regime Identification}

After estimation, the script extracts smoothed state probabilities:
\[
\Pr(s_t=j \mid \mathcal{F}_T),
\]
where \(\mathcal{F}_T\) denotes the full sample information set.

The high-risk regime is defined as the regime with the higher
probability-weighted CARDI mean:
\begin{equation}
j^\star
=
\arg\max_{j \in \{1,2\}}
\frac{\sum_t \Pr(s_t=j \mid \mathcal{F}_T) y_t}
{\sum_t \Pr(s_t=j \mid \mathcal{F}_T)}.
\end{equation}

The high-risk probability is:
\begin{equation}
\pi_t^{H}
=
\Pr(s_t=j^\star \mid \mathcal{F}_T).
\end{equation}

\subsection{Regime Switch Dates}

For a threshold \(c \in \{0.5,0.8\}\), a switch into the high-risk regime is
defined by:
\begin{equation}
\pi_t^H \geq c
\quad \text{and} \quad
\pi_{t-1}^H < c.
\end{equation}

Each switch date is matched to the nearest individual event and nearest event
cluster. The script reports both calendar-day and trading-day distances.

\section{Output Files}

\begin{longtable}{ll}
\toprule
Output & Description\\
\midrule
\texttt{event\_clusters.xlsx} & Individual and clustered event metadata\\
\texttt{chow\_known\_breakpoint\_tests.xlsx} & Known-breakpoint Chow test results\\
\texttt{intervention\_newey\_west\_tests.xlsx} & Interrupted time-series regressions\\
\texttt{reference\_event\_dummy\_newey\_west.xlsx} & Legacy event-dummy regressions with HAC SEs\\
\texttt{markov\_switching\_results.xlsx} & Regime probabilities and switch-event matches\\
\texttt{FRM\_Event.csv} & Smoothed CARDI series and log differences\\
\texttt{RegEvent\_SmoothFRMDiff*.csv} & Reference event-dummy regression outputs\\
\texttt{Plot\_Event.png} & Legacy CARDI event plot\\
\texttt{MarkovSwitching\_CARDI\_*.png} & Markov regime probability plots\\
\bottomrule
\end{longtable}

\section{Interpretation}

The tests answer related but distinct questions. The Chow test asks whether the
local regression relationship changes around a known event date. The
interrupted time-series regressions decompose the event association into an
immediate pulse, a persistent level shift, and a post-event trend change. The
reference event-dummy model asks whether recent event occurrence predicts
changes in smoothed CARDI. The Markov-switching model asks whether CARDI enters
a high-risk latent regime and whether such regime transitions occur near known
events.

\section*{References}

\begin{thebibliography}{9}

\bibitem[Bernal, Cummins, and Gasparrini(2017)]{bernal2017}
Bernal, J. L., Cummins, S., and Gasparrini, A. (2017).
\newblock Interrupted time series regression for the evaluation of public
health interventions: a tutorial.
\newblock \emph{International Journal of Epidemiology}, 46(1), 348--355.
\newblock \url{https://doi.org/10.1093/ije/dyw098}

\bibitem[Chow(1960)]{chow1960}
Chow, G. C. (1960).
\newblock Tests of equality between sets of coefficients in two linear
regressions.
\newblock \emph{Econometrica}, 28(3), 591--605.
\newblock \url{https://doi.org/10.2307/1910133}

\bibitem[Hamilton(1989)]{hamilton1989}
Hamilton, J. D. (1989).
\newblock A new approach to the economic analysis of nonstationary time series
and the business cycle.
\newblock \emph{Econometrica}, 57(2), 357--384.
\newblock \url{https://doi.org/10.2307/1912559}

\bibitem[Newey and West(1987)]{neweywest1987}
Newey, W. K., and West, K. D. (1987).
\newblock A simple, positive semi-definite, heteroskedasticity and
autocorrelation consistent covariance matrix.
\newblock \emph{Econometrica}, 55(3), 703--708.
\newblock \url{https://doi.org/10.2307/1913610}

\end{thebibliography}

\end{document}
